% Options for packages loaded elsewhere
\PassOptionsToPackage{unicode}{hyperref}
\PassOptionsToPackage{hyphens}{url}
\PassOptionsToPackage{dvipsnames,svgnames,x11names}{xcolor}
%
\documentclass[
  10pt,
]{article}
\usepackage{amsmath,amssymb}
\usepackage{iftex}
\ifPDFTeX
  \usepackage[T1]{fontenc}
  \usepackage[utf8]{inputenc}
  \usepackage{textcomp} % provide euro and other symbols
\else % if luatex or xetex
  \usepackage{unicode-math} % this also loads fontspec
  \defaultfontfeatures{Scale=MatchLowercase}
  \defaultfontfeatures[\rmfamily]{Ligatures=TeX,Scale=1}
\fi
\usepackage{lmodern}
\ifPDFTeX\else
  % xetex/luatex font selection
\fi
% Use upquote if available, for straight quotes in verbatim environments
\IfFileExists{upquote.sty}{\usepackage{upquote}}{}
\IfFileExists{microtype.sty}{% use microtype if available
  \usepackage[]{microtype}
  \UseMicrotypeSet[protrusion]{basicmath} % disable protrusion for tt fonts
}{}
\makeatletter
\@ifundefined{KOMAClassName}{% if non-KOMA class
  \IfFileExists{parskip.sty}{%
    \usepackage{parskip}
  }{% else
    \setlength{\parindent}{0pt}
    \setlength{\parskip}{6pt plus 2pt minus 1pt}}
}{% if KOMA class
  \KOMAoptions{parskip=half}}
\makeatother
\usepackage{xcolor}
\usepackage[margin=0.72in]{geometry}
\usepackage{longtable,booktabs,array}
\usepackage{calc} % for calculating minipage widths
% Correct order of tables after \paragraph or \subparagraph
\usepackage{etoolbox}
\makeatletter
\patchcmd\longtable{\par}{\if@noskipsec\mbox{}\fi\par}{}{}
\makeatother
% Allow footnotes in longtable head/foot
\IfFileExists{footnotehyper.sty}{\usepackage{footnotehyper}}{\usepackage{footnote}}
\makesavenoteenv{longtable}
\setlength{\emergencystretch}{3em} % prevent overfull lines
\providecommand{\tightlist}{%
  \setlength{\itemsep}{0pt}\setlength{\parskip}{0pt}}
\setcounter{secnumdepth}{5}
\usepackage{amsmath,amssymb,mathtools}
\usepackage{booktabs,longtable,array}
\usepackage{microtype}
\setlength{\parindent}{0pt}
\setlength{\parskip}{5pt plus 1pt minus 1pt}
\renewcommand{\arraystretch}{1.13}
\ifLuaTeX
  \usepackage{selnolig}  % disable illegal ligatures
\fi
\usepackage{bookmark}
\IfFileExists{xurl.sty}{\usepackage{xurl}}{} % add URL line breaks if available
\urlstyle{same}
\hypersetup{
  pdftitle={Production-Theory Closure of the PT/BFM-Constrained Three-Loop 3PI Pilot},
  pdfauthor={Angus Muffatti},
  colorlinks=true,
  linkcolor={blue},
  filecolor={Maroon},
  citecolor={Blue},
  urlcolor={blue},
  pdfcreator={LaTeX via pandoc}}

\title{Production-Theory Closure of the PT/BFM-Constrained Three-Loop
3PI Pilot}
\usepackage{etoolbox}
\makeatletter
\providecommand{\subtitle}[1]{% add subtitle to \maketitle
  \apptocmd{\@title}{\par {\large #1 \par}}{}{}
}
\makeatother
\subtitle{Exact topology ledger, complete thermal tensor spaces,
counterterm closure basis, analytic benchmark hierarchy, and
preregistered observable contract}
\author{Angus Muffatti}
\date{24 August 2026}

\begin{document}
\maketitle

{
\hypersetup{linkcolor=}
\setcounter{tocdepth}{3}
\tableofcontents
}
\section{Executive verdict}\label{executive-verdict}

The five residual theory tasks identified after v1.8 are complete:

\begin{enumerate}
\def\labelenumi{\arabic{enumi}.}
\tightlist
\item
  the exact generic three-loop 3PI diagram ledger is frozen;
\item
  the finite-temperature tensor spaces are complete for the declared
  symmetry class;
\item
  every local counterterm, composite insertion and initial-surface
  signature required by power counting is mapped;
\item
  the analytic and reduced-numerical benchmark hierarchy is frozen;
\item
  the physical-observable and error contract is preregistered.
\end{enumerate}

\[
\boxed{12/12\ \text{preflight checks passed}},
\qquad
\boxed{10/10\ \text{unit tests passed}}.
\]

\[
\boxed{\text{PASS FOR UNIT TEST AND PILOT}}
\]

This is not an unconditional production pass. A finite non-Abelian 3PI
truncation is not guaranteed to preserve gauge identities or possess a
conventional all-order renormalization proof. The pilot tests those
properties through Ward, Slavnov-Taylor, Nielsen, cutoff, KMS,
conservation and gauge-singlet spectral gates.

The v1.8 package already supplied the matched retarded benchmark, seed
STI closure, singlet ladder and executable solver specification. The
present work removes the remaining ambiguity in what equations, tensor
components, counterterms, limits and observables the pilot means.

\section{Frozen scope}\label{frozen-scope}

The high-temperature electroweak-symmetric portal is

\[
\mathcal L_Y=-y_D\,\overline Q_L H D_R+\mathrm{h.c.},
\]

coupled to \(SU(3)_c\times SU(2)_L\times U(1)_Y\), with vectorlike
colored singlet \(D\). The temperature anchor is
\(T=1.002\times10^8\,\mathrm{GeV}\).

Dynamical two-point functions:

\[
G_H,S_Q,S_D,D_g,D_W,D_B,G_{c_3},G_{c_2}.
\]

Dynamical three-point vertices include portal, matter-gauge,
scalar-gauge, three-gauge and gauge-ghost families. Four-point vertices
remain at renormalized classical values at this truncation order, with
their counterterms retained.

\section{Exact three-loop 3PI
functional}\label{exact-three-loop-3pi-functional}

\[
\Gamma_2=\Gamma_2^0+\Gamma_2^{\rm int}.
\]

\[
\begin{aligned}
\Gamma_2^0={}&-\frac18D^2V_4^0+\frac{i}{6}D^3V_3V_3^0-iD\Delta^2UU_0
+\frac{i}{24}D^4V_4V_4^0+\frac18D^5V_3^2V_4^0,\\
\Gamma_2^{\rm int}={}&-\frac{i}{12}D^3V_3^2+\frac{i}{2}D\Delta^2U^2-\frac{i}{48}D^4V_4^2
-\frac{i}{24}D^6V_3^4+\frac{i}{3}D^3\Delta^3U^3V_3+\frac{i}{4}D^2\Delta^4U^4.
\end{aligned}
\]

All contour signs and multiplicities are carried by superfield tensors.
The ledger contains \textbf{11} unique generic topologies; every row
passes \(L=I-V+1\).

\subsection{Why the ledger matters}\label{why-the-ledger-matters}

Self-energies, dynamical vertices and the singlet Bethe-Salpeter kernel
must be differentiated from one functional. Independent coding risks
double counting and breaks the conserving relation

\[
K=\frac{\delta\Sigma_H}{\delta G_H}.
\]

\subsection{Frozen topology summary}\label{frozen-topology-summary}

\begin{longtable}[]{@{}
  >{\raggedright\arraybackslash}p{(\columnwidth - 6\tabcolsep) * \real{0.2500}}
  >{\raggedright\arraybackslash}p{(\columnwidth - 6\tabcolsep) * \real{0.2500}}
  >{\raggedright\arraybackslash}p{(\columnwidth - 6\tabcolsep) * \real{0.2500}}
  >{\raggedright\arraybackslash}p{(\columnwidth - 6\tabcolsep) * \real{0.2500}}@{}}
\toprule\noalign{}
\begin{minipage}[b]{\linewidth}\raggedright
ID
\end{minipage} & \begin{minipage}[b]{\linewidth}\raggedright
L
\end{minipage} & \begin{minipage}[b]{\linewidth}\raggedright
Coefficient
\end{minipage} & \begin{minipage}[b]{\linewidth}\raggedright
Topology
\end{minipage} \\
\midrule\noalign{}
\endhead
\bottomrule\noalign{}
\endlastfoot
G20\_B4\_BARE & 2 & -1/8 & bosonic double-bubble/seagull \\
G20\_B33\_MIX & 2 & +i/6 & bosonic sunset with one dressed and one bare
cubic vertex \\
G20\_F3\_MIX & 2 & -i & Grassmann sunset with one dressed and one bare
boson-Grassmann vertex \\
G30\_B44\_MIX & 3 & +i/24 & bosonic basketball with one dressed and one
bare quartic vertex \\
G30\_B334 & 3 & +1/8 & bosonic squint with two dressed cubics and one
bare quartic \\
G2I\_B33 & 2 & -i/12 & bosonic sunset with two dressed cubic vertices \\
G2I\_F3 & 2 & +i/2 & Grassmann sunset with two dressed vertices \\
G3I\_B44 & 3 & -i/48 & bosonic basketball with two dressed quartics \\
G3I\_B3333 & 3 & -i/24 & bosonic tetrahedron/Mercedes with four dressed
cubic vertices \\
G3I\_FFFB & 3 & +i/3 & Grassmann triangle joined to a dressed bosonic
cubic vertex \\
G3I\_FFFF & 3 & +i/4 & Grassmann box with four dressed boson-Grassmann
vertices \\
\end{longtable}

\section{Complete tensor spaces}\label{complete-tensor-spaces}

The state is homogeneous, CP-even and parity-even, with thermal
four-velocity \(u^\mu\). Angular dependence is carried by spherical
harmonics.

The complete Clifford basis has rank 16. A generic fermion-gauge vertex
has

\[
64=16_L+48_T
\]

components. The STI fixes the longitudinal matrix sector; all 48
transverse components are dynamical. The measured transversality
residual is \(1.556e-17\).

Other allocations:

\begin{itemize}
\tightlist
\item
  scalar/ghost-gauge vertex: one longitudinal plus three transverse
  components;
\item
  chiral portal vertex: four projected matrices per orientation;
\item
  matter-ghost kernel: all 16 Clifford components;
\item
  three-gauge vertex: 27 fully transverse component tensors before
  Bose/color/permutation reduction.
\end{itemize}

The three-gauge transversality residual is \(4.445e-16\). These are
storage component spaces, not counts of final independent physical
scalar form factors.

\section{Renormalization closure}\label{renormalization-closure}

For the declared renormalizable theory and initial-state class, every
local structure allowed by power counting and symmetries has a mapped
counterterm or initial kernel. The superficial degree is

\[
\omega=4-d_I-n_B-n_c-\frac32n_F.
\]

The matrix contains \textbf{32} signatures and no missing entries.

The basis includes two-point, split three-point, Yukawa, quartic,
truncation-longitudinal, composite-operator and initial-surface
structures. For \(\mathcal O_H=H^\dagger H\), it includes operator
renormalization, mixing and a local two-source contact term. The initial
state includes \(\alpha_2,\alpha_3,\alpha_4\) prepared through an
interacting imaginary-time leg plus a UV-soft finite deformation.

No free state-independent counterterm exists for an unaccompanied
\(p<6\) harmonic. Persistence after selector/state charges vanish is a
hard failure.

This is a power-counting and symmetry closure basis, not an all-order
theorem for finite non-Abelian 3PI. Cutoff independence during evolution
is the real test.

\section{Benchmark hierarchy}\label{benchmark-hierarchy}

The frozen hierarchy covers free unequal-time normalization, equilibrium
KMS, Abelian Ward, pure Yukawa, linear response, kinetic/AMY, narrow
width, factorization-scale cancellation and conserving singlet BSE
limits.

\begin{longtable}[]{@{}ll@{}}
\toprule\noalign{}
Check & Result \\
\midrule\noalign{}
\endhead
\bottomrule\noalign{}
\endlastfoot
Free equal-time spectral value & 0.000e+00 \\
Free commutator derivative error & 8.817e-13 \\
KMS residual & 1.474e-16 \\
Abelian Ward residual & 4.487e-16 \\
Finite-memory/Markov error & 7.919e-04 \\
Narrow-width area error & 7.639e-03 \\
\end{longtable}

All current analytic checks pass.

\section{Observable and error
contract}\label{observable-and-error-contract}

Claim-bearing quantities are limited to:

\begin{itemize}
\tightlist
\item
  Nielsen-stable complex poles;
\item
  integrated matched portal rate;
\item
  gauge-singlet \(H^\dagger H\) spectral density;
\item
  conserved sector energy and gauge charges;
\item
  downstream \(T_5/T_0\).
\end{itemize}

Conventional gauge-fixed off-shell elementary line shapes and individual
ghost/longitudinal dressings are diagnostics.

Preregistered systematic scans cover momentum and angular resolution,
timestep, memory, gauge parameter, separation scale, renormalization
scale, tensor-basis size and truncation controls. Signed shifts and
covariance are reported; broken identities are not error bars.

\section{Pilot configuration}\label{pilot-configuration}

\[
N_r=96,\qquad\ell_{\max}=4,\qquad N_t=4096,\qquad N_{\rm mem}=256.
\]

\[
\xi\in\{0,0.5,1,2\},\qquad q_*/T\in\{0.15,0.25,0.40\}.
\]

The package uses deterministic reductions in validation, HDF5
checkpoints and full double/complex-double precision.

\section{Validation}\label{validation}

\begin{longtable}[]{@{}ll@{}}
\toprule\noalign{}
Gate & Status \\
\midrule\noalign{}
\endhead
\bottomrule\noalign{}
\endlastfoot
factorization\_scale & PASS \\
on\_shell\_anchor & PASS \\
STI\_seed & PASS \\
singlet\_BSE & PASS \\
singlet\_positivity & PASS \\
KMS\_noise\_positivity & PASS \\
diagram\_ledger & PASS \\
fermion\_vertex\_basis & PASS \\
three\_gauge\_basis & PASS \\
counterterm\_closure\_basis & PASS \\
analytic\_benchmarks & PASS \\
observable\_contract & PASS \\
\end{longtable}

All 12 preflight gates pass. The Python suite contains 10 tests; all
pass without warnings.

\section{Production authorization}\label{production-authorization}

Unit and eight-GPU pilot runs are authorized. Production remains blocked
until the pilot preserves simultaneously

\[
\boxed{\text{Ward/ST}+\text{Nielsen}+\text{KMS}+\text{conservation}+\text{cutoff/}q_*\text{ independence}+\text{singlet positivity}}.
\]

The pilot's central scientific question is whether the finite three-loop
3PI truncation maintains physical gauge consistency over long real-time
evolution. No remaining small analytic exercise can settle that
honestly.

\section{Final status}\label{final-status}

\begin{longtable}[]{@{}
  >{\raggedright\arraybackslash}p{(\columnwidth - 2\tabcolsep) * \real{0.5000}}
  >{\raggedright\arraybackslash}p{(\columnwidth - 2\tabcolsep) * \real{0.5000}}@{}}
\toprule\noalign{}
\begin{minipage}[b]{\linewidth}\raggedright
Item
\end{minipage} & \begin{minipage}[b]{\linewidth}\raggedright
Status
\end{minipage} \\
\midrule\noalign{}
\endhead
\bottomrule\noalign{}
\endlastfoot
Exact three-loop 3PI ledger & PASS \\
Complete component tensor spaces & PASS \\
Counterterm/initial-state closure basis & PASS FOR DECLARED
TRUNCATION \\
Analytic benchmark hierarchy & PASS \\
Observable/error contract & PASS \\
Automated preflight & 12/12 PASS \\
Unit tests & 10/10 PASS \\
Pilot authorization & APPROVED \\
Production authorization & CONDITIONAL ON PILOT \\
All-order non-Abelian 3PI renormalization theorem & NOT CLAIMED \\
\end{longtable}

The remaining uncertainty now sits inside the evolving non-Abelian
correlator/vertex system, where it is measured by identities and
physical controls rather than hidden in an ansatz.

\section{References}\label{references}

\begin{enumerate}
\def\labelenumi{\arabic{enumi}.}
\tightlist
\item
  J. Berges, \emph{n-Particle irreducible effective action techniques
  for gauge theories}, arXiv:hep-ph/0401172.
\item
  J. Berges, \emph{Introduction to Nonequilibrium Quantum Field Theory},
  arXiv:hep-ph/0409233.
\item
  U. Reinosa and J. Serreau, \emph{2PI effective action for gauge
  theories: Renormalization}, arXiv:hep-th/0605023.
\item
  M. Garny and M. M. Muller, \emph{Kadanoff-Baym Equations with
  Non-Gaussian Initial Conditions: The Equilibrium Limit},
  arXiv:0904.3600.
\item
  S. Borsanyi and U. Reinosa, \emph{Renormalized nonequilibrium quantum
  field theory: Scalar fields}, arXiv:0809.0496.
\item
  Project v1.8, \emph{Pre-HPC Closure of the Gauge-Covariant q-D-H
  Portal}.
\end{enumerate}

\section{Machine-readable
deliverables}\label{machine-readable-deliverables}

The package includes the topology ledger, vertex/tensor catalogues,
counterterm matrix, solver specifications, observable contract,
preflight driver, tests, launch checklist, acceptance matrix and
checksums.

\end{document}
